% !TEX root = scientific-computing.tex

\chapter{Probability Distributions and Pseudorandom Numbers}
\label{ch:random}

One big part of scientific computation is the subject of Monte Carlo
simulations in which random numbers are used to model some situation.
We'll spend the next part of this course covering this subject.


\section{Probability}\label{probability}

To fully understand how to use random numbers and how Monte Carlo
simulations work, we need some basic probability. When understanding
probability there are two types continuous and discrete, which for what
we need will come down to using floating point or integers.

\subsection{Discrete Probability}

We'll first discuss discrete probability. A \emph{discrete probability
distribution} is a finite set (often will be a subset of integers.) For
example, let's use the set
\begin{align*}
A=\{1,2,3,4,5,6\}
\end{align*}
%
and an event, $X$ is subset of these numbers. For example, let
$X=\{1\}$, then the probability that this event occurs in the ratio of
the number of elements in each set or \[P(X)=\frac{N(X)}{N(A)}\]

In this case, $P(X)=1/6$ and that means that the probability that the
number 1 comes up is 1/6. Think about this in the case of rolling a die.
This says that the probability that 1 comes up is 1/6.

\subsection{Continuous Probability}

Another type of probability distribution is called a continuous
distribution and in this case we'll only consider a type of continuous
probability distribution called a uniform distribution.

Let's consider a set $A=\{x \; | \; 0 \leq x \leq 1\}$ or all real
numbers between 0 and 1. Events are still subsets of the set $A$,
however the probability that events occur is the fraction of the set.

For example, if $X=\{x \; | \; 1/3 \leq x \leq 1/2\}$, then $P(X)$
is the fraction
\[P(X) = \frac{\frac{1}{2}-\frac{1}{3}}{1-0}=\frac{1}{6}\]


\section{Pseudo Random Number Generator}

\href{https://en.wikipedia.org/wiki/Pseudorandom_number_generator}{The
Wikipedia Pseudo-random Number Generator page} gives an overview of the
subject and a lot of technical details. In short, a truly random number
on a computer is very difficult to generate and generally not necessary
because a pseudorandom number is sufficient. I also won't try to explain
this in general term since you'd need a significant mathematical
background.

A pseudo-random number generator is a function that produces a sequence
of numbers that act like random numbers. Let's examine what this means
in terms of the discrete probability with set $A=\{1,2,3,4,5,6\}$.

If a pseudo-random number generator produces a sequence from this set
then the sequence should have the following property:

\begin{itemize}

\item
  If the event $X=\{i\}$ for any $i$ between 1 and 6, then
  $P(X)\approx 1/6$. And by approximately, as the sequence gets
  larger, the approximation becomes closer to 1/6.
\end{itemize}

Is this enough? No, the sequence \[\{1,2,3,4,5,6,1,2,3,4,5,6,\ldots\}\]

satisfies the above property, but I don't think anyone would consider
this random. Another property would be:

\begin{itemize}

\item
  If we know the sequence $\{a_1,a_2,a_3,\ldots,a_n\}$ then we can't
  predict the next number $a_{n+1}$.
\end{itemize}

This is obviously violated in the sequence above.

A little more technical definition of a sequence of pseudo-random
numbers Let\\ $(a_1,a_2,a_3, \ldots)$ be a random sequence. Typically we
mean the following properties need to hold:

\begin{enumerate}
\def\labelenumi{\arabic{enumi}.}

\item
  any number in the range 1 to 6 is equally likely to occur.
\item
  Take N random numbers and let $s_n$ be the number of times the
  number $n$ occurs. The fraction $s_n/N$ should go to 1/6 in the
  limit as $n\rightarrow \infty$.
\item
  Knowing the sequence $s_1, s_2, \ldots, s_k$ does not allow us to
  predict $s_{k+1}$
\end{enumerate}

If instead we use floating point numbers, there are a few different
properties. Assume that the floating point number is in the range
$0 < x < 1$. Then the sequence $(a_1,a_2,a_3, \ldots)$ is random
if

\begin{enumerate}
\def\labelenumi{\arabic{enumi}.}
\setcounter{enumi}{2}

\item
  Let $s_{[c,d]}(a_1,a_2,a_3, \ldots)$ be the total of the numbers
  that satisfy $c < a_k < d$. Then as the number of random numbers
  approach infinity $s_{[c,d]}=d-c$
\end{enumerate}

\hypertarget{using-julia-to-simulate-the-rolling-of-a-die}{%
\subsection{Using Julia to simulate the rolling of a
die}\label{using-julia-to-simulate-the-rolling-of-a-die}}

First, the main commands that are built-in to Julia are listed in the
\href{https://docs.julialang.org/en/latest/stdlib/Random/#Random-Numbers-1}{Julia
Manual for Random Numbers}. We can generate 100 random numbers between 1
and 6 using
%
\begin{jlblock}
S=rand(1:6,100)
\end{jlblock}
%
and notice that if you rerun the command, you'll get a different
sequence of random numbers. We can check that this is doing what we
expect by checking the probability that we get a 1 (or any other
number).

To check the count of the number 1 that appears, try

\begin{verbatim}
count(a->a==1,S)
\end{verbatim}

\subsection{Exercise}

Convert this to a probability. Is it close to what you expect? Try
changing the number of random numbers used to larger numbers. Does your
answer get closer to what you expect?

\hypertarget{floating-point-random-numbers}{%
\subsection{Floating Point random
numbers}\label{floating-point-random-numbers}}

Let's look a floating point random numbers between 0 and 1. To generate
a sequence (array) of such numbers, type

\begin{verbatim}
S=rand(100)
\end{verbatim}

and we can check the number of values less than 0.1 with the command:

\begin{verbatim}
count(a->a<0.1,S)
\end{verbatim}

and if we want the number of values between 0.4 and 0.7 then type

\begin{verbatim}
count(a->0.4<a<0.7,S)
\end{verbatim}

\subsection{Exercise}

Change the number of random numbers used to much larger set and return
the above commands. What are the results? Is this what you expect?

How to do find the fraction of number that are less that 0.25 and
greater than 0.75?


\hypertarget{calculating-pi-using-pseudo-random-numbers}{%
\section{\texorpdfstring{Calculating $\pi$ using pseudo random
numbers}{Calculating \textbackslash pi using pseudo random numbers}}\label{calculating-pi-using-pseudo-random-numbers}}

\begin{itemize}
\item
  \href{https://en.wikipedia.org/wiki/Buffon\%27s_needle}{Buffon's
  Needle Experiment (18th C.)}
\item
  Circle in the Square.

  Consider the square
  $\{ (x,y)\; | \; 0\leq x\leq 1, 0 \leq y \leq 1\}$ and the quarter
  circle that falls within $x^{2}+y^{2} \leq 1$. The area within the
  circle is the fraction of the cirle in the square or $\pi/4$.

  We can use this to estimate a value of $\pi$.

  \begin{enumerate}
  \item  Randomly choose points $x$ and $y$ in the range [0,1].
  \item  Count all points that fall within the circle.
  \item  Find the fraction of the points in the circle.
  \item  Since this fraction should be $\pi/4$, multiply by 4 to estimate $\pi$.
  \end{enumerate}

  This can be done in julia using

  \begin{enumerate}
  \item \jlb{S1=rand(100,2)}
  \item \jlb{S2 = mapslices(x-> x[1]^2+x[2]^2,S1;dims=[2])}
  \item \jlb{S3 = count(a -> a < 1,S2)}
  \item \jlb{fr=S3/100}
  \item \jlb{est=4*fr}
  \end{enumerate}
which has the result \jlc{print(est)}. And all of these can be combined into one command:

\begin{jlblock}
4*count(a->a<1,mapslices(x->x[1]^2+x[2]^2,rand(100,2); dims=[2]))
\end{jlblock}
\end{itemize}

\subsection{Exercise}
\begin{itemize}
\item  Write a function that take a positive integer $n$, in and performs
  the approximate-$\pi$ calculation described above using $n$
  points. It should return the approximate value of $\pi$.
\item  Test it will high values of $n$ and time it.
\item  for $N=10^5$, $N=10^6$ and $N=10^7$ find the relative error of
  the estimate using your function and using the built-in value
  \texttt{pi}.
\end{itemize}

